\documentclass[11pt,a4paper]{article}
\usepackage[T1]{fontenc}
\usepackage{lmodern,amsmath,amssymb,booktabs,microtype}
\usepackage[margin=25mm]{geometry}
\usepackage[hidelinks]{hyperref}
\newcommand{\interp}{\par\smallskip\noindent\textit{Interpretation.}\ }
\newcommand{\R}{\mathbb{R}}
\title{\small Session 2 --- Climate block\\[0.6em]
\LARGE Climate Dynamics as an Initial-Value Problem}
\author{Gauthier Vermandel\\\small \'Ecole polytechnique}
\date{}
\begin{document}
\maketitle

This note describes how an emissions path generates carbon concentrations,
radiative forcing and temperature changes. The model is deliberately small:
three carbon reservoirs, one methane concentration and two temperatures.
We first pose a continuous-time ordinary differential equation (ODE) system,
then derive the explicit time-stepping scheme used in
\texttt{Climate\_Models.py}. Every stock, flow and initial condition has a
distinct role.

\section{What is given, what is solved?}
Unlike the economic block, the stand-alone climate block contains no
optimization. Emissions are prescribed inputs. Once parameters, initial
states and an emissions scenario are specified, the task is to solve an
initial-value problem, also called a Cauchy problem.

Time $t$ is measured in years on a finite interval $[t_0,t_f]$.
Let $M_a,M_u,M_\ell$ denote atmospheric, upper-ocean and lower-ocean carbon
stocks in gigatonnes of carbon (GtC). Let $q$ be atmospheric methane
concentration in parts per billion (ppb), and $T_a,T_\ell$ temperature
anomalies relative to the model's reference climate in degrees Celsius.
The state and input vectors are
\begin{equation}\label{eq:vectors}
 x=(M_a,M_u,M_\ell,q,T_a,T_\ell)^\top,\qquad
 e=(e_{\mathrm{CO_2}},e_{\mathrm{CH_4}})^\top .
\end{equation}
\interp The six states summarize the system's inherited physical condition.
The prescribed inputs are a carbon-dioxide mass flow
$e_{\mathrm{CO_2}}$ in GtCO$_2$/year and a methane concentration source
$e_{\mathrm{CH_4}}$ in ppb/year. A mass-based methane series must therefore
be converted before it can enter this model.

For fixed parameters $\theta$, the complete mathematical task is
\begin{equation}\label{eq:cauchy}
 \boxed{\dot x(t)=f(x(t),e(t);\theta),\qquad
 x(t_0)=x^0,\qquad t\in[t_0,t_f].}
\end{equation}
\interp The derivative is the instantaneous change of the state per year.
Six first-order equations require six initial values. No terminal
temperature target or terminal state is imposed: future climate is the
output. An inverse problem (find emissions meeting a target) or a control
problem (choose emissions to optimize welfare) would be a different task.

All components of $x^0$ must be specified at the same calendar date.
Carbon stocks are positive, methane concentration is positive, and
temperature anomalies may have either sign. Stocks cannot be replaced
by emissions, nor can an anomaly be interpreted as an absolute temperature.

\section{Physical equations and their interpretation}
The model is triangular: emissions affect carbon and methane; these
determine radiative forcing; forcing drives temperatures. There is no
temperature feedback on the carbon cycle in this particular benchmark.
That omission is a modeling restriction, not a general physical claim.

\paragraph{Carbon exchange.}
Let $(m_a,m_u,m_\ell)$ be reference carbon stocks, with each entry positive.
The annual transfer coefficients are $b_{12},b_{23}\geq0$.
Define reverse rates and the CO$_2$-to-carbon conversion by
\begin{equation}\label{eq:rates}
 b_{21}=b_{12}\frac{m_a}{m_u},\qquad
 b_{32}=b_{23}\frac{m_u}{m_\ell},\qquad
 \xi=\frac{12}{44}=\frac{3}{11}.
\end{equation}
\interp Reverse rates balance the exchange flows at the reference stocks.
The factor $\xi$ converts a mass of CO$_2$ into its carbon content;
it is not the fraction of emissions that remains in the atmosphere.

Writing $M=(M_a,M_u,M_\ell)^\top$ and $v_a=(1,0,0)^\top$, carbon dynamics are
\begin{equation}\label{eq:carbon}
 \dot M=BM+\xi e_{\mathrm{CO_2}}v_a,\qquad
 B=\begin{pmatrix}
 -b_{12}&b_{21}&0\\
 b_{12}&-(b_{21}+b_{23})&b_{32}\\
 0&b_{23}&-b_{32}
 \end{pmatrix}.
\end{equation}
\interp Each diagonal entry is an outflow rate from a reservoir; each
off-diagonal term transfers carbon from another reservoir. Emissions enter
the atmosphere. There is no direct atmosphere--lower-ocean exchange.
The upper ocean links the fast atmospheric reservoir to the slow deep one.

Two identities are useful for understanding and checking the system:
\begin{equation}\label{eq:conservation}
 B\begin{pmatrix}m_a\\m_u\\m_\ell\end{pmatrix}=0,\qquad
 \frac{d}{dt}(M_a+M_u+M_\ell)=\xi e_{\mathrm{CO_2}}.
\end{equation}
\interp Without emissions, the reference stocks are stationary, and
internal transfers conserve total carbon. Zero emissions do not generally
restore the reference climate: an initially elevated total stock remains
inside these three modeled reservoirs and is redistributed among them.
The zero eigenvalue associated with carbon conservation is not a numerical
instability.

\paragraph{Methane.}
Let $\bar q>0$ be the reference concentration and $\delta_q>0$ the removal
rate. The one-box methane equation is
\begin{equation}\label{eq:methane}
 \dot q=-\delta_q(q-\bar q)+e_{\mathrm{CH_4}},
 \qquad \tau_q=\delta_q^{-1}.
\end{equation}
\interp Excess methane decays with characteristic lifetime $\tau_q$,
while the prescribed source replenishes it. The restoring term includes
an implicit background source maintaining $\bar q$; setting the additional
input to zero returns methane toward $\bar q$, not toward zero.

\paragraph{Radiative forcing.}
Let $F_{2\times}>0$ denote forcing from doubled atmospheric CO$_2$,
and $\alpha_q\geq0$ the methane forcing coefficient. The forcing diagnostic is
\begin{equation}\label{eq:forcing}
 F(M_a,q)=F_{2\times}\frac{\log(M_a/m_a)}{\log2}
       +\alpha_q(\sqrt q-\sqrt{\bar q}).
\end{equation}
\interp Forcing, in W/m$^2$, is the radiative energy imbalance induced by
the gas concentrations. Doubling $M_a$ from its reference value contributes
$F_{2\times}$; the methane term measures forcing relative to $\bar q$.
Forcing is an algebraic output, not an additional state requiring its
own initial condition. This simplified formula omits gas-overlap effects,
aerosols and other external forcing components.

\paragraph{Temperature response.}
Let $S>0$ be equilibrium climate sensitivity, and set
$\lambda=F_{2\times}/S$. With $c_1,c_4>0$ and $c_3>0$,
\begin{equation}\label{eq:temperatures}
 \begin{aligned}
 \dot T_a&=c_1\big[F-\lambda T_a-c_3(T_a-T_\ell)\big],\\
 \dot T_\ell&=c_4(T_a-T_\ell).
 \end{aligned}
\end{equation}
\interp Surface temperature rises with forcing, falls with radiative
restoration $\lambda T_a$, and loses heat to a colder deep ocean.
Deep-ocean temperature adjusts more slowly toward the surface temperature.
Here $\lambda$ and $c_3$ have units W/m$^2$/K, $c_1$ has units
K/(year W/m$^2$), and $c_4$ has units year$^{-1}$.

For constant forcing and an equilibrium in both temperature boxes,
\begin{equation}\label{eq:equilibrium}
 T_a^*=T_\ell^*=\frac{F}{\lambda},
 \qquad F=F_{2\times}\ \Longrightarrow\ T_a^*=S.
\end{equation}
\interp Equilibrium climate sensitivity is the eventual response to
sustained doubling forcing after ocean adjustment. It is not the
temperature reached immediately, nor the warming at a fixed near-term
date under a changing emissions scenario.

\section{Well-posedness and discretization}
Writing an ODE is not enough: its domain and assumptions must allow a
solution, and the numerical scheme must respect the properties we intend
to study. We separate the continuous problem from its approximation.

\paragraph{Existence, uniqueness and dependence on inputs.}
Suppose the emissions inputs are bounded and piecewise continuous on
$[t_0,t_f]$, reference stocks are strictly positive, and the initial carbon
and methane states are strictly positive. On the domain where $M_a>0$
and $q>0$, the right-hand side is locally Lipschitz in the state.
The linear reservoir equations and the two-box temperature equations
therefore give a unique solution on any finite interval that stays in
this domain; solutions depend continuously on initial data and inputs.

For nonnegative emissions, the carbon matrix has nonnegative off-diagonal
entries and the methane background source is positive. These properties
preserve nonnegative stocks; strictly positive initial stocks stay positive
on finite intervals. Bounded inputs then keep stocks finite and the forcing
well-defined, so the solution extends through $t_f$. If negative emissions
are allowed, positivity must be checked: an input that exhausts atmospheric
carbon takes the logarithm outside the model's domain.

\paragraph{Forward Euler on a calendar grid.}
Let $t_n=t_0+n\Delta$, with $n=0,\ldots,N$. Holding inputs constant
over each interval gives the explicit approximation
\begin{equation}\label{eq:euler}
 x_{n+1}=x_n+\Delta f(x_n,e_n;\theta),\qquad x_0=x^0.
\end{equation}
\interp A derivative is a rate per year; multiplying by $\Delta$ converts
it into a change over the interval. Forward Euler evaluates every
right-hand side at the old state. Under usual smoothness assumptions its
global error is first order in $\Delta$; halving the step is a basic
convergence check, not a change in the physical model.

In particular, the carbon and methane steps are
\begin{equation}\label{eq:gassteps}
 \begin{aligned}
 M_{n+1}&=(I+\Delta B)M_n+\Delta\xi e_{\mathrm{CO_2},n}v_a,\\
 q_{n+1}&=\bar q+(1-\Delta\delta_q)(q_n-\bar q)
                         +\Delta e_{\mathrm{CH_4},n}.
 \end{aligned}
\end{equation}
\interp If emissions are stored as annual rates, both source terms require
$\Delta$. If instead the input is already a period total, it must not be
multiplied by $\Delta$ again. The revised stand-alone module consistently
uses annual rates for both inputs.

Sufficient conditions for nonnegative reservoir weights are
\begin{equation}\label{eq:positive}
 \Delta\max\{b_{12},\,b_{21}+b_{23},\,b_{32}\}\leq1,
 \qquad \Delta\delta_q\leq1.
\end{equation}
\interp A discrete step should not remove more of a reservoir than it
contains. These conditions prevent negative weights in the explicit gas
updates; they are sufficient conditions, not necessary ones for every
particular trajectory.

For temperatures, write $T=(T_a,T_\ell)^\top$. The update is
\begin{equation}\label{eq:tempstep}
 T_{n+1}=(I+\Delta D)T_n+
      \Delta\begin{pmatrix}c_1\\0\end{pmatrix}F(M_{a,n},q_n),
 \quad
 D=\begin{pmatrix}
 -c_1(\lambda+c_3)&c_1c_3\\
 c_4&-c_4
 \end{pmatrix}.
\end{equation}
\interp The forcing and exchange terms are evaluated at the beginning of
the interval. For a stable homogeneous discrete temperature response,
each eigenvalue $\eta$ of $D$ must satisfy $|1+\Delta\eta|<1$.
Thus a stable continuous ODE can still have an unstable explicit
discretization if the time step is too large.

\section{From the equations to the notebook}
The notebook first considers a stationary reference climate, then feeds
historical emissions into the model and compares temperature anomalies.
Future scenarios are conditional simulations, not forecasts independent
of their emissions assumptions.

\paragraph{Code mapping.}
The names in \texttt{Params} and in the simulation matrix correspond to
the mathematical objects as follows.
\begin{center}\small
\begin{tabular}{@{}lll@{}}
\toprule
Mathematical object & Python name & Unit\\
\midrule
$M_a,M_u,M_\ell$ & \texttt{M\_AT, M\_UP, M\_LO} & GtC\\
$q$ & \texttt{M\_CH4} & ppb (not a mass)\\
$T_a,T_\ell$ & \texttt{T\_AT, T\_LO} & $^\circ$C anomaly\\
$e_{\mathrm{CO_2}}$ & \texttt{E} & GtCO$_2$/year\\
$e_{\mathrm{CH_4}}$ & \texttt{E\_CH4} & ppb/year\\
$F$ & \texttt{F} & W/m$^2$\\
$m_a,m_u,m_\ell$ & \texttt{mat, mup, mlo} & GtC\\
$S,F_{2\times}$ & \texttt{T2XCO2, F2XCO2} & $^\circ$C; W/m$^2$\\
$\bar q,\delta_q,\alpha_q$ &
 \texttt{mch4, delta\_ch4, alpha\_ch4} & As in \eqref{eq:methane}--\eqref{eq:forcing}\\
\bottomrule
\end{tabular}
\end{center}

Row zero stores the initial climate state. To compute row $j\geq1$,
\texttt{update\_path} reads states and emissions from row $j-1$.
The requested final date is included. Consequently, an emissions value
recorded at year 2020 acts over the interval starting in 2020; it is not
the emissions value for the preceding interval.

The forcing stored in \texttt{F[j]} is evaluated from the states in row
$j-1$ and used for that interval's temperature update. It therefore
represents interval forcing, whereas \texttt{T\_AT[j]} is the end-of-interval
temperature. The initial forcing entry can remain \texttt{NaN}: it is
not an independent state. Emissions at the initial date, in contrast,
must be provided because they drive the first transition.

\paragraph{Calibration and initial conditions.}
The default simulation uses $\Delta=5$ years and begins in 2015.
The historical exercise overrides the step to one year and resets the
states to a preindustrial reference. Those are different initial-value
problems and must not be combined accidentally.
\begin{center}\small
\begin{tabular}{@{}lll@{}}
\toprule
Quantity & Default value & Role\\
\midrule
$(m_a,m_u,m_\ell)$ & $(588,360,1720)$ & Reference carbon stocks\\
$(M_a^0,M_u^0,M_\ell^0)$ & $(851,460,1740)$ & Initial carbon stocks\\
$(b_{12},b_{23})$ & $(0.024,0.0014)$ & Annual transfer rates\\
$(\bar q,q^0,\delta_q)$ & $(722,1860,1/9)$ & Methane reference, initial value, decay\\
$(F_{2\times},S)$ & $(3.3613,3.1)$ & Forcing scale and climate sensitivity\\
$(c_1,c_3,c_4)$ & $(0.0201,0.0176,0.005)$ & Temperature-response coefficients\\
$(T_a^0,T_\ell^0)$ & $(0.85,0.007)$ & Initial temperature anomalies\\
$\alpha_q$ & $0.036$ & Methane forcing coefficient\\
\bottomrule
\end{tabular}
\end{center}

\paragraph{Data are not model states.}
The historical CO$_2$ series is an emissions flow: convert tonnes
to gigatonnes before applying $\xi$. Do not feed CO$_2$-equivalent methane
emissions directly into a concentration equation: the warming-potential
conversion and the conversion from methane mass to ppb are separate
assumptions. Match all series by calendar year rather than array position.
For temperature comparisons, harmonize reference periods for anomalies;
otherwise a level difference may be a baseline mismatch rather than
a model error. These simplified exercises do not estimate all model
parameters or reproduce a full climate-model calibration.

\paragraph{Reproducible checks.}
Start at $(m_a,m_u,m_\ell,\bar q,0,0)$ with zero inputs: every state should
remain stationary. Next add a CO$_2$ pulse and verify that the change in
total carbon equals $\Delta\xi e_{\mathrm{CO_2}}$ at each step.
Then run a constant-forcing temperature experiment and check convergence
to \eqref{eq:equilibrium}. Finally compare annual and finer-step simulations
at common dates, preserving the same annual emissions inputs.
These tests distinguish equation errors, unit errors and discretization
errors before interpreting differences between scenarios.

\end{document}
